\documentclass[11pt]{article}
\usepackage[margin=1in]{geometry}
\usepackage{times}
\usepackage[hidelinks]{hyperref}
\usepackage{titlesec}
\titlespacing*{\section}{0pt}{6pt}{2pt}
\titleformat{\section}{\normalfont\large\bfseries}{}{0pt}{}
\setlength{\parskip}{4pt}
\setlength{\parindent}{0pt}

\begin{document}

\begin{center}
{\Large\bfseries Measuring and Forecasting the Impact of the 2025 U.S.\
Government Shutdown on New York Airport Delays}\\[4pt]
{\normalsize Student Investigator: Pranish Uprety, Department of Computer
Science, Caldwell University}\\
{\normalsize Faculty Advisor: Dr.\ Sushil Poudel}\\
{\normalsize ICUNJ Research Symposium --- Description of Research}
\end{center}

\section*{Rationale for the Proposed Research}
Flight delays are a familiar frustration for travelers, but they are also a
major economic problem: on average, a single minute of delay costs an airline
about \$100 once crew, fuel, and aircraft costs are counted. Delays are rarely
isolated. Because airports share the same aircraft, crews, and airspace, trouble
at one busy airport spreads to its neighbors, much like a traffic jam that backs
up onto connecting roads. In the fall of 2025, the United States experienced its
longest-ever federal government shutdown (43 days). Air-traffic controllers were
required to work without pay, staffing shortages grew, and the Federal Aviation
Administration eventually ordered airlines to cut flights at 40 major airports.
This unusual, well-documented event is a rare ``natural experiment'': it lets us
measure what happens to the air-travel system when a key government service is
disrupted. This research asks three simple but important questions about the New
York City airports (Newark, JFK, and LaGuardia), the busiest and most
interconnected airport cluster in the country: (1) How much worse did delays and
cancellations actually get? (2) Was the increase caused by the shutdown, or could
it be explained away by bad weather? (3) Could the disruption have been predicted
ahead of time using only information the public already had?

\section*{Relevant Background Literature}
Three strands of prior work motivate this study. First, foundational network
research by Fleurquin and colleagues showed that delays ``ripple'' across the
U.S.\ airport network in a measurable, systemic way, establishing that airports
cannot be studied in isolation~[1]. Second, recent computer-science work uses
graph-based machine learning to predict delays across airport networks; for
example, Cai and colleagues built a model that separates the geographic and
operational connections between airports to improve delay prediction~[2], and
Zhu and colleagues developed ``CausalNet,'' which tries to learn cause-and-effect
relationships between airports rather than assuming a fixed map~[3]. These models
predict or diagnose delays but do not measure the effect of a specific real-world
disruption. Third, an economic study from the Federal Reserve Bank of Richmond
analyzed 76 million flights and estimated that the 2025 shutdown added roughly
one million extra hours of delay nationwide (about \$66 million), even after
accounting for weather~[4]. That study works at the national level; the present
project complements it by zooming into a single, tightly coupled metro region and
by adding a forecasting test. The industry cost figure used throughout comes from
Airlines for America, which reports a direct operating cost of \$100.76 per
minute of delay for 2024~[5].

\section*{Research Procedures}
The project uses only free, publicly available data, so it is fully reproducible.
The main data source is the U.S.\ Bureau of Transportation Statistics on-time
performance database, which records the scheduled and actual departure time of
essentially every U.S.\ flight. Hourly weather observations for the three New
York airports are taken from the Iowa Environmental Mesonet archive. The
procedure has four steps. (1)~\textbf{Preparation:} clean the flight records and
organize them into an ``airport-by-hour'' table showing the average delay at each
airport in each hour, for the shutdown period (October--November 2025) and the
same months of the previous, normal year (2024) for comparison. (2)~\textbf{
Measurement:} compare the two years directly, including how often flights were
cancelled outright. (3)~\textbf{Attribution:} estimate the ``excess'' delay
caused specifically by the shutdown using three independent methods (a simple
prior-year comparison and two machine-learning models, one of which also accounts
for weather), then check whether the three methods agree. (4)~\textbf{
Forecasting:} test whether a model that knows \emph{only} the calendar, the
scheduled number of flights, and the public fact that a shutdown was underway,
with no access to actual delays, could anticipate the disruption better than a
normal-year expectation.

\section*{Specific Means of Evaluation of Findings}
Each finding is judged against a clear, quantitative standard rather than
opinion. For the size of the delay increase, statistical significance tests
(a Welch $t$-test and a bootstrap confidence interval) determine whether the
change is real or could be normal year-to-year variation. For the cost estimate,
credibility comes from \emph{agreement} among the three independent methods: if
they converge on a similar number, the estimate is trustworthy, and the
weather-adjusted method specifically tests whether weather, not the shutdown,
was responsible. For the forecasting question, the shutdown-aware model is scored
against a naive ``normal-year'' baseline using standard accuracy measures
($R^2$ and mean absolute error) on data the model never saw during training, so
it cannot simply memorize the answer.

\section*{Significance of the Potential Findings}
This work turns a major national event into concrete, local, dollar-denominated
evidence of what happens to air travel when government services are interrupted.
A clear, weather-controlled measurement of the shutdown's cost to the New York
airports would give policymakers and the public a tangible sense of the stakes,
and complement existing national estimates. Showing that the disruption was
partly \emph{forecastable} from public information suggests that similar events
could be anticipated and managed rather than merely reacted to. More broadly, the
results lay the groundwork for future ``intervention-aware'' models that could
help decide where and when to slow traffic to minimize nationwide disruption.

\section*{Preliminary Work and Feasibility}
This proposal builds on substantial work already completed during summer 2026,
so the timetable is realistic. The full data pipeline is built and the analyses
have been run. Preliminary results are strong and consistent: average departure
delay rose 107\% at JFK and 149\% at LaGuardia (both highly statistically
significant), cancellations rose three- to eight-fold, and the three independent
methods agree that the shutdown caused roughly 0.40--0.42 million extra minutes
of delay, about \$41 million in direct airline cost, which weather does not
explain. The forecasting test succeeds: a model using only advance information
outperforms a normal-year baseline at all three airports. A draft research paper
documenting these results is already in preparation. The remaining work,
finalizing figures, writing, and revising for presentation, is well within the
symposium timeline.

\section*{References}
{\small
[1] P.\ Fleurquin, J.\ J.\ Ramasco, and V.\ M.\ Egu\'iluz, ``Systemic delay
propagation in the US airport network,'' \textit{Scientific Reports}, vol.\ 3,
p.\ 1159, 2013.

[2] K.\ Cai, Y.\ Li, Y.-P.\ Fang, and Y.\ Zhu, ``A geographical and operational
deep graph convolutional approach for flight delay prediction,'' \textit{Chinese
Journal of Aeronautics}, vol.\ 36, no.\ 3, pp.\ 357--367, 2023.
                          
[3] Q.\ Zhu, S.\ Chen, T.\ Guo, Y.\ Lv, and W.\ Du, ``CausalNet: A
Granger-causality-driven spatio-temporal graph neural network for flight delay
propagation,'' arXiv:2407.15185, 2024.

[4] R.\ T.-N.\ Wong, ``The Value of Government Services: Evidence from Flight
Delays During the 2025 Shutdown,'' Federal Reserve Bank of Richmond, Economic
Brief 25-45, 2025.

[5] Airlines for America, ``U.S.\ Passenger Carrier Delay Costs,'' 2025.
CY2024 direct operating cost of delay: \$100.76 per block-minute.
}

\end{document}
